% Souffler dans une solution de bleu de bromothymol (exercice 7, chapitre 1)
\begin{tikzpicture}[x=1cm,y=1cm,
  etat/.style={font=\otLight\fontsize{7.5}{9.2}\selectfont,text=otInk,align=center,text width=23mm,anchor=north}]
\pic at (0,0) {becher={largeur=1.6,hauteur=1.9,niveau=1.15,liquide=pcVert}};
\pic at (0.15,0.22) {paille={hauteur=2.7}};
\draw[pcRepere] (0.21,2.5) -- (0.95,2.75) node[pcLegende,anchor=west]{paille};
\draw[pcRepere] (-0.45,0.6) -- (-1.15,0.95) node[pcLegende,anchor=east]{eau + BBT};
\node[etat] at (0,-0.3) {L'élève souffle dans la solution.};
\draw[pcFleche] (1.5,1) -- (2.5,1);
\pic at (3.6,0) {tube={niveau=1.15,liquide=pcVert}};
\pic at (6.1,0) {tube={niveau=1.15,liquide=pcJaune}};
\pic at (8.6,0) {tube={niveau=1.15,liquide=pcBleu}};
\node[etat] at (3.6,-0.3) {avant};
\node[etat] at (6.1,-0.3) {après avoir soufflé};
\node[etat] at (8.6,-0.3) {après ajout de solution d'hydroxyde de sodium};
\end{tikzpicture}
\par\smallskip\legende{Teintes successives de la solution de bleu de bromothymol.}
